\documentclass[12pt,a4paper]{ctexart}
\usepackage[margin=2.5cm]{geometry}
\usepackage{hyperref}

\hypersetup{colorlinks=true,linkcolor=black,urlcolor=blue}
\setlength{\parindent}{2em}
\setlength{\parskip}{0.35em}

\newcommand{\code}[1]{\nolinkurl{#1}}
\newenvironment{dialogue}
  {\begin{quote}\small\setlength{\parskip}{0.45em}}
  {\end{quote}}

\title{并行程序设计 Lab4：AI 使用报告}
\author{刘蔚霖\quad 2412727}
\date{2026 年 6 月 9 日}

\begin{document}
\maketitle

\section{关键对话}

\begin{dialogue}
\textbf{我：}阅读 Lab4 要求和口令猜测代码，规划 MPI 加速方案。Backup 保持不动，
并分别实现基础版和进阶版。

\textbf{AI：}基础版使用阻塞式 Master-Worker。Rank 0 维护唯一 PT 优先队列，将完整 PT
动态分配给空闲 Worker；Worker 在本地生成候选口令并执行 MD5。进阶版使用生成线程、
有界缓存和非阻塞 MPI，将候选生成、传输与加密组织为流水线。

\end{dialogue}

\begin{dialogue}
\textbf{我：}Base 版本具体如何满足 Master-Worker 要求？

\textbf{AI：}Rank 0 作为 Master，独占概率优先队列并调用 \code{PopNextTask()}
取得概率最高的预终结式（Pre-Terminal，PT）。Master 将完整 PT 封装成
\code{WorkUnit} 后发送给 Worker。Worker 调用 \code{GenerateRange()} 在本地生成
候选口令，并执行 NEON MD5，完成后只返回处理数量和计时统计。
\end{dialogue}

\begin{dialogue}
\textbf{我：}Base 为什么采用完整 PT 作为任务，而不发送生成后的候选口令？

\textbf{AI：}PT 描述比完整候选字符串批次紧凑。发送 PT 后，Worker 可以利用本地训练好的
PCFG 模型独立派生候选口令，节点间只传输任务描述和统计结果。该数据分配方式减少网络传输量，
并将候选生成与 MD5 计算同时分散到多个 Worker。
\end{dialogue}

\begin{dialogue}
\textbf{我：}Base 如何处理不同 PT 工作量差异造成的负载不均衡？

\textbf{AI：}Master 使用 \code{MPI\_ANY\_SOURCE} 接收任意 Worker 的完成结果。
最先完成的 Worker 会立即获得下一个 PT，因此执行较短任务的 Worker 可以连续领取工作。
该策略属于动态任务调度，能够缓解完整 PT 大小不同造成的静态分配不均衡。
\end{dialogue}

\begin{dialogue}
\textbf{我：}\code{WorkUnit} 为什么需要单独设计？

\textbf{AI：}\code{WorkUnit} 同时保存 PT 和最后一个 segment 的半开区间
\code{[begin,end)}。Base 使用覆盖完整 PT 的范围；Advanced 可以将大型 PT 切分为多个
互不重叠的范围任务。统一的数据结构使两个版本共享范围生成接口，并保证切分任务之间没有重复或遗漏。
\end{dialogue}

\begin{dialogue}
\textbf{我：}PT 和字符串含有 STL 动态数据，MPI 应该怎样发送？

\textbf{AI：}不能直接发送 C++ 对象的内存布局。使用 \code{MPI\_Pack} 按固定顺序写入
segment 数量、类型、长度、索引和范围边界，接收端再以相同顺序调用 \code{MPI\_Unpack}。
这种字段级序列化传输逻辑内容，不依赖对象中的指针地址。
\end{dialogue}

\begin{dialogue}
\textbf{我：}Advanced 版本的实现思路是什么？

\textbf{AI：}Advanced 将流水线分为候选生成、MPI 调度和加密三个部分。Rank 0 启动多个
Pthread 生成线程，每个线程拥有独立任务队列和有界 READY 缓存。生成线程将
\code{WorkUnit} 展开成口令批次；Rank 0 主线程轮询 READY 缓存并发送批次；其他 Rank
作为加密 Worker，接收批次并执行 NEON MD5。
\end{dialogue}

\begin{dialogue}
\textbf{我：}为什么生成线程不能直接调用 MPI？

\textbf{AI：}程序通过 \code{MPI\_Init\_thread} 请求
\code{MPI\_THREAD\_FUNNELED}。该线程级别规定只有初始化 MPI 的主线程调用 MPI 接口。
生成线程只负责候选生成和缓存操作，Rank 0 主线程统一处理发送、接收和请求完成，
从而保持明确的 MPI 调用边界。
\end{dialogue}

\begin{dialogue}
\textbf{我：}Advanced 如何给多个生成线程分配任务？

\textbf{AI：}Rank 0 从优先队列连续取得 PT，并使用 \code{SplitTask()} 将大型 PT
切分成范围任务。程序根据 \code{GuessCount()} 估计每个 \code{WorkUnit} 的候选数量，
将下一任务分配给预计累计负载最小的生成线程。该策略以候选数量近似生成工作量。
\end{dialogue}

\begin{dialogue}
\textbf{我：}Advanced 为什么使用非阻塞 MPI？

\textbf{AI：}Rank 0 需要同时维护多个加密 Worker，并让通信与候选生成、MD5 计算重叠。
主线程使用 \code{MPI\_Isend} 发送批次，使用 \code{MPI\_Irecv} 接收统计结果，
再通过 \code{MPI\_Testany} 和 \code{MPI\_Waitany} 处理最先完成的 Worker。
每个 Worker 使用独立发送缓冲区，缓冲区在对应请求完成前保持有效。
\end{dialogue}

\begin{dialogue}
\textbf{我：}为什么 Advanced 设计更复杂，但服务器结果可能慢于 Base？

\textbf{AI：}Base 发送紧凑 PT 描述，Worker 在本地生成并哈希。Advanced 发送已经生成的
完整候选字符串，每条口令都需要打包长度和字符内容。非阻塞通信能够重叠部分等待时间，
但无法消除序列化和数据传输成本。当通信时间占墙钟时间的比例较高时，增加加密 Worker
只能缩短哈希阶段，整体性能会进入通信饱和区间。
\end{dialogue}

\begin{dialogue}
\textbf{我：}为什么增加生成线程不一定继续加速？

\textbf{AI：}增加生成线程只缩短候选生成阶段。当生成端已经能够持续提供 READY 批次后，
墙钟时间主要由 Rank 0 通信和加密 Worker 决定。额外生成线程还会与 MPI 主线程争用处理器
和内存带宽，因此生成关键路径可能继续下降，而整体墙钟时间保持不变或上升。
\end{dialogue}

\begin{dialogue}
\textbf{我：}批次越大是否一定越快？

\textbf{AI：}较大批次可以减少消息次数，但会延长批次发布前的组装等待，并增加尾部任务不均衡。
较小批次能够更早发送，使生成、通信和加密保持细粒度重叠。最佳批次大小取决于消息启动开销、
序列化成本和 Worker 处理时间之间的平衡。
\end{dialogue}

\begin{dialogue}
\textbf{我：}程序如何保证所有 Worker 正常退出？

\textbf{AI：}Master 在任务耗尽后向每个 Worker 发送独立停止消息。Worker 根据消息标签区分
工作任务与停止任务，退出接收循环后进入 \code{MPI\_Finalize}。Advanced 主线程还会等待
所有未完成发送和结果接收请求结束，再发送停止消息并释放线程和缓存。
\end{dialogue}

\begin{dialogue}
\textbf{我：}Base 输出显示处理数量超过 \code{--generate}，是否有问题？

\textbf{AI：}程序按完整 PT 或完整范围任务处理，并在累计数量越过目标后停止，因此最终处理数量
可能略高于目标值。不同进程数和参数配置得到一致的正式处理数量，说明动态调度、任务切分和
MPI 消息没有造成任务重复或遗漏。
\end{dialogue}

\end{document}
